\documentclass[10pt,a4paper]{amsart}
\usepackage{amsmath,amssymb,amsthm,graphicx}
% Retained minimal preamble for the Stacks Project chapter more-algebra.tex (More on Algebra).
% The chapter file pinned-original-source.tex carries only \input{preamble} on its first line, so no
% preamble text lives inside the chapter itself.  The Stacks Project uses one single shared preamble for
% every chapter; the snapshot stages it identically as sources/stacks-chow/preamble.tex and
% sources/stacks-etale-cohomology/preamble.tex.  The block below is copied mechanically, line for line,
% from sources/stacks-chow/preamble.tex lines 14-18 (the Xy-pic setup) and 154-217 (the native theorem
% declarations together with the mathematical macros).  Nothing was added, removed or redefined; the
% wrapper class and the amsmath/amssymb/amsthm/graphicx base packages are supplied by the assembler.
% Chapter: more-algebra.tex -- More on Algebra.
% For commutative diagrams we use Xy-pic
\usepackage[all]{xy}

% We use 2cell for 2-commutative diagrams.
\xyoption{2cell}

% Theorem environments.
%
\theoremstyle{plain}
\newtheorem{theorem}[subsection]{Theorem}
\newtheorem{proposition}[subsection]{Proposition}
\newtheorem{lemma}[subsection]{Lemma}

\theoremstyle{definition}
\newtheorem{definition}[subsection]{Definition}
\newtheorem{example}[subsection]{Example}
\newtheorem{exercise}[subsection]{Exercise}
\newtheorem{situation}[subsection]{Situation}

\theoremstyle{remark}
\newtheorem{remark}[subsection]{Remark}
\newtheorem{remarks}[subsection]{Remarks}

\numberwithin{equation}{subsection}

% Macros
%
\def\lim{\mathop{\mathrm{lim}}\nolimits}
\def\colim{\mathop{\mathrm{colim}}\nolimits}
\def\Spec{\mathop{\mathrm{Spec}}}
\def\Hom{\mathop{\mathrm{Hom}}\nolimits}
\def\Ext{\mathop{\mathrm{Ext}}\nolimits}
\def\SheafHom{\mathop{\mathcal{H}\!\mathit{om}}\nolimits}
\def\SheafExt{\mathop{\mathcal{E}\!\mathit{xt}}\nolimits}
\def\Sch{\mathit{Sch}}
\def\Mor{\mathop{\mathrm{Mor}}\nolimits}
\def\Ob{\mathop{\mathrm{Ob}}\nolimits}
\def\Sh{\mathop{\mathit{Sh}}\nolimits}
\def\NL{\mathop{N\!L}\nolimits}
\def\CH{\mathop{\mathrm{CH}}\nolimits}
\def\proetale{{pro\text{-}\acute{e}tale}}
\def\etale{{\acute{e}tale}}
\def\QCoh{\mathit{QCoh}}
\def\Ker{\mathop{\mathrm{Ker}}}
\def\Im{\mathop{\mathrm{Im}}}
\def\Coker{\mathop{\mathrm{Coker}}}
\def\Coim{\mathop{\mathrm{Coim}}}

% Boxtimes
%
\DeclareMathSymbol{\boxtimes}{\mathbin}{AMSa}{"02}

%
% Macros for moduli stacks/spaces
%
\def\QCohstack{\mathcal{QC}\!\mathit{oh}}
\def\Cohstack{\mathcal{C}\!\mathit{oh}}
\def\Spacesstack{\mathcal{S}\!\mathit{paces}}
\def\Quotfunctor{\mathrm{Quot}}
\def\Hilbfunctor{\mathrm{Hilb}}
\def\Curvesstack{\mathcal{C}\!\mathit{urves}}
\def\Polarizedstack{\mathcal{P}\!\mathit{olarized}}
\def\Complexesstack{\mathcal{C}\!\mathit{omplexes}}
% \Pic is the operator that assigns to X its picard group, usage \Pic(X)
% \Picardstack_{X/B} denotes the Picard stack of X over B
% \Picardfunctor_{X/B} denotes the Picard functor of X over B
\def\Pic{\mathop{\mathrm{Pic}}\nolimits}
\def\Picardstack{\mathcal{P}\!\mathit{ic}}
\def\Picardfunctor{\mathrm{Pic}}
\def\Deformationcategory{\mathcal{D}\!\mathit{ef}}
\usepackage{cleveref}
\crefname{originallemma}{original lemma}{original lemmas}
\Crefname{originallemma}{Original lemma}{Original lemmas}
\makeatletter
\newlabel{algebra-lemma-resolution-by-finite-free}{{00LP}{source}{}{sourcepointer-1}{}}
\newlabel{algebra-lemma-resolution-by-finite-free@cref}{{[originallemma][0][]00LP}{[1][0][]source}{}{}{}}
\makeatother
\crefname{originallemma}{original lemma}{original lemmas}
\Crefname{originallemma}{Original lemma}{Original lemmas}
\makeatletter
\newlabel{algebra-lemma-finitely-presented-localization-free}{{0GWM}{source}{}{sourcepointer-2}{}}
\newlabel{algebra-lemma-finitely-presented-localization-free@cref}{{[originallemma][0][]0GWM}{[1][0][]source}{}{}{}}
\makeatother
\crefname{originallemma}{original lemma}{original lemmas}
\Crefname{originallemma}{Original lemma}{Original lemmas}
\makeatletter
\newlabel{algebra-lemma-finite-projective}{{00NX}{source}{}{sourcepointer-3}{}}
\newlabel{algebra-lemma-finite-projective@cref}{{[originallemma][0][]00NX}{[1][0][]source}{}{}{}}
\makeatother
\crefname{originalproposition}{original proposition}{original propositions}
\Crefname{originalproposition}{Original proposition}{Original propositions}
\makeatletter
\newlabel{algebra-proposition-regular-finite-gl-dim}{{00O7}{source}{}{sourcepointer-4}{}}
\newlabel{algebra-proposition-regular-finite-gl-dim@cref}{{[originalproposition][0][]00O7}{[1][0][]source}{}{}{}}
\makeatother
\crefname{originallemma}{original lemma}{original lemmas}
\Crefname{originallemma}{Original lemma}{Original lemmas}
\makeatletter
\newlabel{algebra-lemma-independent-resolution}{{00O5}{source}{}{sourcepointer-5}{}}
\newlabel{algebra-lemma-independent-resolution@cref}{{[originallemma][0][]00O5}{[1][0][]source}{}{}{}}
\makeatother
\crefname{originallemma}{original lemma}{original lemmas}
\Crefname{originallemma}{Original lemma}{Original lemmas}
\makeatletter
\newlabel{algebra-lemma-Noetherian-topology}{{00FQ}{source}{}{sourcepointer-6}{}}
\newlabel{algebra-lemma-Noetherian-topology@cref}{{[originallemma][0][]00FQ}{[1][0][]source}{}{}{}}
\makeatother
\crefname{originallemma}{original lemma}{original lemmas}
\Crefname{originallemma}{Original lemma}{Original lemmas}
\makeatletter
\newlabel{lemma-cohomology-perfect}{{066U}{source}{}{sourcepointer-7}{}}
\newlabel{lemma-cohomology-perfect@cref}{{[originallemma][0][]066U}{[1][0][]source}{}{}{}}
\makeatother
\title{Over a regular ring perfect modules are exactly the finite modules and perfect complexes are those with bounded finite cohomology}
\author{The Stacks Project Authors}
\date{stacks-project:more-algebra.tex:local-snapshot-20261012; GNU FDL 1.2 or later}
\begin{document}
\maketitle
\setcounter{subsection}{1}
\noindent\textit{Editorial note.} All delivered mathematical text is the author's own, copied from the pinned Stacks Project chapter file more-algebra.tex; no statement or proof step was rewritten, reordered or removed. The counted statement is the labeled lemma environment at source lines 19114-19122 (canonical tag 066Z) and the counted proof is the proof environment at source lines 19124-19173, which immediately follows it and is its own proof.\par\medskip
\noindent\textit{Editorial note.} The proof refers six times to labelled results of another chapter (resolutions by finite free modules, openness of the free locus, finite projectives over local rings, finite global dimension of regular local rings, the syzygy consequence, and Noetherianity of the spectrum) and once to a criterion of this chapter; all seven are carried as navigation stubs with their roles recorded in the reasons.\par\medskip
\noindent\textit{Editorial note.} The chapter file pinned-original-source.tex carries only an input of the shared preamble on its first line, so no preamble text lives inside the chapter. The retained file sources/stacks-more-algebra/retained-preamble.tex is a mechanical, line-for-line copy of the project's single shared preamble as staged in the snapshot (Xy-pic setup plus the native theorem declarations and mathematical macros). Nothing was added, removed or redefined, and the excerpt serves only as preamble, never as a span coordinate.\par\medskip
\noindent\textit{Editorial note.} The source URL is the source-level Stacks Project URL, as required when more than one card is produced from one source; the per-claim tag is carried by claim\_key and by this note. The counted statement and proof cite no bibliography entry, so this card needs no bibliography.\par\medskip
\noindent\textit{Allowed standard prerequisites (editorial).}\begin{itemize}
\item A finite module over a Noetherian ring admits a resolution by finite free modules, and the set of primes at which the module is free is open. used at proof lines 19126-19134. Produces the resolution whose syzygies are studied and makes the loci U\_i open subsets of the spectrum.
\item A finite projective module over a local ring is free, and a finitely generated module of finite projective dimension at a prime can be replaced by a free one after passing to a large syzygy. used at proof lines 19135-19149. Shows that the loci U\_i are increasing and that a large syzygy is free at every prime, hence projective.
\item The spectrum of a Noetherian ring is a Noetherian topological space, so an increasing sequence of open subsets that covers it stabilizes, and for a bounded complex with finite cohomology each cohomology module is perfect. used at proof lines 19150-19154 and 19163-19173. Provides the finite index n at which the covering stabilizes and reduces part (2) to part (1) applied to the cohomology modules.
\end{itemize}\medskip
\ifdefined\hypertarget\hypertarget{sourcepointer-1}{}\fi
\noindent\textit{Source reference.} Original lemma 00LP: sources/stacks-chow/tags.txt, source lines 795--795. The reference supplies the existence of a resolution by finite free modules for a finite module, used at proof lines 19126-19131; the covering argument that follows is inside the counted span.\par\medskip
\ifdefined\hypertarget\hypertarget{sourcepointer-2}{}\fi
\noindent\textit{Source reference.} Original lemma 0GWM: sources/stacks-chow/tags.txt, source lines 21077--21077. The reference makes the locus of primes at which a syzygy is free an open subset, used at proof lines 19132-19134; the monotonicity and covering arguments for those loci are inside the counted span.\par\medskip
\ifdefined\hypertarget\hypertarget{sourcepointer-3}{}\fi
\noindent\textit{Source reference.} Original lemma 00NX: sources/stacks-chow/tags.txt, source lines 875--875. The reference records that a finite projective module over a local ring is free, used at proof lines 19139-19140 and again at proof lines 19153-19154; the sequence of inclusions and the stabilization are inside the counted span.\par\medskip
\ifdefined\hypertarget\hypertarget{sourcepointer-4}{}\fi
\noindent\textit{Source reference.} Original proposition 00O7: sources/stacks-chow/tags.txt, source lines 885--885. The reference gives finite global dimension of the regular local rings, used at proof lines 19143-19145 as the reason the syzygies eventually become free at each prime; everything done with that fact is inside the counted span.\par\medskip
\ifdefined\hypertarget\hypertarget{sourcepointer-5}{}\fi
\noindent\textit{Source reference.} Original lemma 00O5: sources/stacks-chow/tags.txt, source lines 883--883. The reference is the standard consequence of finite global dimension that makes a large syzygy finite projective, used at proof lines 19146-19149; the rest of the covering argument is inside the counted span.\par\medskip
\ifdefined\hypertarget\hypertarget{sourcepointer-6}{}\fi
\noindent\textit{Source reference.} Original lemma 00FQ: sources/stacks-chow/tags.txt, source lines 580--580. The reference records that the spectrum of a Noetherian ring is Noetherian as a topological space, used at proof lines 19150-19152 to conclude that the increasing cover by the loci stabilizes; the conclusion drawn from stabilization is inside the counted span.\par\medskip
\ifdefined\hypertarget\hypertarget{sourcepointer-7}{}\fi
\noindent\textit{Source reference.} Original lemma 066U: sources/stacks-more-algebra/pinned-original-source.tex, source lines 18987--18987. The reference is the criterion applied at proof lines 19170-19172 to deduce perfectness of a bounded complex from perfectness of its cohomology modules, which is part (1) proved inside the counted span.\par\medskip
\begin{lemma}
\label{lemma-regular-perfect}
Let $R$ be a regular ring. Then
\begin{enumerate}
\item an $R$-module is perfect if and only if it is a finite $R$-module, and
\item a complex of $R$-modules $K^\bullet$ is perfect if and only
if $K^\bullet \in D^b(R)$ and each $H^i(K^\bullet)$ is a finite $R$-module.
\end{enumerate}
\end{lemma}

\begin{proof}
Any perfect $R$-module is finite by definition. Conversely, let $M$ be
a finite $R$-module. Choose a resolution
$$
\ldots \to F_2 \xrightarrow{d_2} F_1 \xrightarrow{d_1} F_0 \to M \to 0
$$
with $F_i$ finite free $R$-modules
(Algebra, Lemma \ref{algebra-lemma-resolution-by-finite-free}).
Set $M_i = \Ker(d_i)$. Denote $U_i \subset \Spec(R)$ the set of primes
$\mathfrak p$ such that $M_{i, \mathfrak p}$ is free; $U_i$ is open by
Algebra, Lemma \ref{algebra-lemma-finitely-presented-localization-free}.
We have a exact sequence $0 \to M_{i + 1} \to F_{i + 1} \to M_i \to 0$.
If $\mathfrak p \in U_i$, then
$0 \to M_{i + 1, \mathfrak p} \to F_{i + 1, \mathfrak p} \to
M_{i, \mathfrak p} \to 0$ splits. Thus $M_{i + 1, \mathfrak p}$
is finite projective, hence free 
(Algebra, Lemma \ref{algebra-lemma-finite-projective}).
This shows that $U_i \subset U_{i + 1}$.
We claim that $\Spec(R) = \bigcup U_i$. Namely, for every prime ideal
$\mathfrak p$ the regular local ring $R_\mathfrak p$ has finite
global dimension by
Algebra, Proposition \ref{algebra-proposition-regular-finite-gl-dim}.
It follows that $M_{i, \mathfrak p}$ is finite projective
(hence free) for
$i \gg 0$ for example by Algebra, Lemma
\ref{algebra-lemma-independent-resolution}.
Since the spectrum of $R$ is Noetherian
(Algebra, Lemma \ref{algebra-lemma-Noetherian-topology})
we conclude that $U_n = \Spec(R)$
for some $n$. Then $M_n$ is a projective $R$-module by
Algebra, Lemma \ref{algebra-lemma-finite-projective}.
Thus
$$
0 \to M_n \to F_n \to \ldots \to F_1 \to M \to 0
$$
is a bounded resolution by finite projective modules
and hence $M$ is perfect.
This proves part (1).

\medskip\noindent
Let $K^\bullet$ be a complex of $R$-modules.
If $K^\bullet$ is perfect, then it is in $D^b(R)$ and it is
quasi-isomorphic to a finite complex of finite projective $R$-modules
so certainly each $H^i(K^\bullet)$ is a finite $R$-module (as $R$ is
Noetherian). Conversely, suppose that $K^\bullet$ is in $D^b(R)$
and each $H^i(K^\bullet)$ is a finite $R$-module. Then by (1) each
$H^i(K^\bullet)$ is a perfect $R$-module, whence $K^\bullet$ is
perfect by
Lemma \ref{lemma-cohomology-perfect}
\end{proof}

\end{document}
